%% kapa: exam=1 round=20 year=2009 period=2 form=A subject=accounting subject_ko=회계학 seq=35 number=35
%% source: Q-Net 감정평가사 기출문제 (cst00309) artlSeq=5024178 file=extracted/1008/5024178/2009년 제20회 감정평가사 1차시험 기출문제/2교시.hwp
%% conversion: kordoc 4.24.51 -> tools/gen1.py (auto); review=proofread-sonnet; flags=none
\begin{question}{35}
(주)한국은 결합원가 ￦420,000을 투입하여 연산품 X, Y, Z를 생산한다. 연산품 X와 Z는 추가가공하여 판매하고 있다. 결합원가 배부방법은 순실현가치법이며, 당기에 생산된 수량은 모두 당기에 판매된다.

{\small\centering
\begin{tabular}{@{}c r r r@{}}
\hline
제품 & \multicolumn{1}{c}{생산수량} & \multicolumn{1}{c}{매출액} & \multicolumn{1}{c}{개별원가(추가가공원가)} \\
\hline
X & 10,000개 & ￦250,000 & ￦100,000 \\
Y & 15,000개 & 200,000 & \\
Z & 20,000개 & 300,000 & ￦150,000 \\
\hline
\end{tabular}\par}
(주)한국은 Y제품을 추가 가공하면 Y제품의 매출액은 ￦550,000이 될 것으로 판단하고 있다. 이 경우에 추가가공원가가 최대 얼마 미만일 때, Y제품을 추가 가공하는 것이 더 유리한가? \par
\begin{choices}
  \item ￦300,000
  \item ￦350,000
  \item ￦400,000
  \item ￦450,000
  \item ￦500,000
\end{choices}
\end{question}
\answer{2}
